\documentclass[11pt,a4paper]{article}
\usepackage[T1]{fontenc}
\usepackage[utf8]{inputenc}
\usepackage[french]{babel}
\usepackage{lmodern}
\usepackage[margin=27mm]{geometry}
\usepackage{xcolor}
\definecolor{navy}{RGB}{23,49,71}
\definecolor{muted}{RGB}{78,89,98}
\newcommand{\editionlink}[2]{\leavevmode\pdfstartlink user{/Subtype /Link /A << /S /URI /URI (#1) >>}\textcolor{navy}{#2}\pdfendlink}
\pagestyle{empty}
\begin{document}
\color{navy}
\noindent\textsc{Bibliothèque de géométrie algébrique}\par
\vspace{22mm}
\noindent{\fontsize{32}{38}\selectfont\bfseries Éléments de géométrie\par algébrique}\par
\vspace{9mm}
\noindent{\Large Alexandre Grothendieck et Jean Dieudonné}\par
\vspace{18mm}
\noindent{\LARGE Édition française}\par
\vspace{3mm}
\noindent{\large Texte corrigé --- lecteur cumulatif EGA I--IV}\par
\vspace{10mm}
\color{muted}
\noindent Cette édition réunit les volumes I, II, III-1, III-2 et IV-1 à IV-4,
y compris les préliminaires (EGA~0) qui leur sont associés.
Les signets et les références internes permettent de parcourir le lecteur.\par
\vspace{5mm}
\noindent Édition corrigée. Le texte diplomatique constitue une édition
distincte. Les notes mathématiques précisent les lectures retenues.\par
\vspace{9mm}
\noindent\editionlink{https://doi.org/10.5281/zenodo.21921588}{Édition française historique : DOI 10.5281/zenodo.21921588}\par
\vspace{3mm}
\noindent\editionlink{https://doi.org/10.5281/zenodo.21921591}{Édition anglaise historique : DOI 10.5281/zenodo.21921591}\par
\vspace{3mm}
\noindent\editionlink{https://github.com/KokunoYumeto/ega-fr}{Sources, fichiers LaTeX et versions : github.com/KokunoYumeto/ega-fr}\par
\vfill
\small
\noindent Préparation des sources et assemblage : OpenAI Codex,
GPT-6 Astra, effort Ultra.\par
\vspace{3mm}
\noindent Le PDF est accompagné du texte LaTeX et du projet éditable,
avec les fichiers nécessaires à la compilation.\par
\end{document}
